\scanpage{238}
\mspage{188}

veux être \unsure{strict}, et donner (si je peux) des démonstrations complètes dans le cadre où je me place à présent, celui des modéliseurs $\mathrm{Cat}$, et en n'utilisant que ce que je sais y démontrer directement à présent, sans me référer à \deleted{\unsure{les}} la théorie $\frac12$ simpliciale ou topologique.

\uline{Théorème} 2 (?) Soit $X$ dans $\mathrm{Cat}$, $n\in\unsure{\mathbb N|\mathcal N}$ un entier $\geqslant 1$. Alors les conditions suivantes sont équivalentes :

$a_n)$ $X$ est $0$-connexe, et \deleted{\unsure{pour}} si $x$ est dans $X$, $\pi_{\unsure{i|n}}(X,x)=0$ pour $i\in[1,n]$ (condition qui ne dépend pas du choix de $x$), i.e. les \inserted{$\Omega^i(X,x)$ pour $i\in[1,n]$ sont $0$-connexes,}

$b_n)$ Pour \deleted{tout} \deleted{\unsure{$F$ faisc. constant}} \deleted{\unsure{faisc.}} groupe \deleted{\unsure{abélien}} $G$, $G\longrightarrow H^0(X,G_X)$ est un iso, et \deleted{\unsure{pour}} $H^1(X,G_X)=0$. De plus si $G$ est abélien, on a aussi $H^i(X,G_X)=0$ pour $i\in[1,n]$.

\uline{Rém.} Pour $n=1$ cette équivalence est $\pm$ tautologique et a été rappelée \inserted{(sous une forme légèrement différente)} au début. (p.~71). On procède par récurrence. On peut donc supposer que \inserted{$n\geqslant 2$, que} $X$ satisfait $a_{n-1},b_{n-1}$, et prouver qu'alors $a_n,b_n$ sont \unsure{équivalentes|équivalents}, i.e. $\pi_n(X,x)=0$ équivaut à $H^n(X,\deleted{\unsure{M}}A)=0$ pour tt gr. ab. $\deleted{\unsure{M}}A$.
% archive sheet number: 238
